% Options for packages loaded elsewhere
\PassOptionsToPackage{unicode}{hyperref}
\PassOptionsToPackage{hyphens}{url}
\documentclass[
  12pt,
]{ctexart}
\usepackage{xcolor}
\usepackage{amsmath,amssymb}
\setcounter{secnumdepth}{-\maxdimen} % remove section numbering
\usepackage{iftex}
\ifPDFTeX
  \usepackage[T1]{fontenc}
  \usepackage[utf8]{inputenc}
  \usepackage{textcomp} % provide euro and other symbols
\else % if luatex or xetex
  \usepackage{unicode-math} % this also loads fontspec
  \defaultfontfeatures{Scale=MatchLowercase}
  \defaultfontfeatures[\rmfamily]{Ligatures=TeX,Scale=1}
\fi
\usepackage{lmodern}
\ifPDFTeX\else
  % xetex/luatex font selection
\fi
% Use upquote if available, for straight quotes in verbatim environments
\IfFileExists{upquote.sty}{\usepackage{upquote}}{}
\IfFileExists{microtype.sty}{% use microtype if available
  \usepackage[]{microtype}
  \UseMicrotypeSet[protrusion]{basicmath} % disable protrusion for tt fonts
}{}
\makeatletter
\@ifundefined{KOMAClassName}{% if non-KOMA class
  \IfFileExists{parskip.sty}{%
    \usepackage{parskip}
  }{% else
    \setlength{\parindent}{0pt}
    \setlength{\parskip}{6pt plus 2pt minus 1pt}}
}{% if KOMA class
  \KOMAoptions{parskip=half}}
\makeatother
\usepackage{longtable,booktabs,array}
\usepackage{calc} % for calculating minipage widths
% Correct order of tables after \paragraph or \subparagraph
\usepackage{etoolbox}
\makeatletter
\patchcmd\longtable{\par}{\if@noskipsec\mbox{}\fi\par}{}{}
\makeatother
% Allow footnotes in longtable head/foot
\IfFileExists{footnotehyper.sty}{\usepackage{footnotehyper}}{\usepackage{footnote}}
\makesavenoteenv{longtable}
\setlength{\emergencystretch}{3em} % prevent overfull lines
\providecommand{\tightlist}{%
  \setlength{\itemsep}{0pt}\setlength{\parskip}{0pt}}
\usepackage{bookmark}
\IfFileExists{xurl.sty}{\usepackage{xurl}}{} % add URL line breaks if available
\urlstyle{same}
\hypersetup{
  hidelinks,
  pdfcreator={LaTeX via pandoc}}

\author{SCX}
\date{2026}

\begin{document}

\section{Definition 1: State-Conditioned Expert
Risk}\label{definition-1-state-conditioned-expert-risk}

\begin{quote}
Definition document for the core mathematical objects of the SCX framework. Defines basic concepts such as states, state-conditioned expert risk, SCX reliability, and data value.
\end{quote}

\begin{center}\rule{0.5\linewidth}{0.5pt}\end{center}

\subsection{1. State Space}\label{state-space}

\subsubsection{Definition 1.1: State}\label{definition-1.1-state}

A state \(s\) is a measurable subset of the input space \(\mathcal{X}\). A state partition
\(\Pi = \{s_1, s_2, \dots, s_K\}\) satisfies:

\begin{enumerate}
\def\labelenumi{\arabic{enumi}.}
\tightlist
\item
  \(\bigcup_{k=1}^K s_k = \mathcal{X}\)
\item
  \(s_k \cap s_{k'} = \emptyset\) for \(k \neq k'\)
\item
  Each \(s_k\) is a measurable set in the Borel \(\sigma\)-algebra of \(\mathcal{X}\)
\end{enumerate}

The mathematical essence of state partitioning is to find a mapping
\(\phi: \mathcal{X} \to \mathcal{S}\) such that:

\[P(Y|X) \approx P(Y|\phi(X))\]

i.e., \(\phi(X)\) is an \textbf{approximate sufficient statistic} for \(Y\) with respect to \(X\).

\subsubsection{Definition 1.2: Soft State Assignment}\label{definition-1.2-soft-state-assignment}

The soft state assignment function \(\gamma: \mathcal{X} \times \mathcal{S} \to [0,1]\)
satisfies:

\[\sum_{s \in \mathcal{S}} \gamma_s(x) = 1, \quad \forall x \in \mathcal{X}\]

Common implementations:

\begin{itemize}
\tightlist
\item
  \textbf{GMM posterior probability}: \(\gamma_s(x) = \frac{\pi_s \mathcal{N}(\phi(x) \mid \mu_s, \Sigma_s)}{\sum_{t \in \mathcal{S}} \pi_t \mathcal{N}(\phi(x) \mid \mu_t, \Sigma_t)}\)
\item
  \textbf{Kernel weighting}: \(\gamma_s(x) = \frac{\kappa(\phi(x), \mu_s)}{\sum_{t \in \mathcal{S}} \kappa(\phi(x), \mu_t)}\)
\item
  \textbf{KNN probability}: \(\gamma_s(x) = \frac{|\text{KNN}(x) \cap s|}{k}\)
\end{itemize}

\begin{center}\rule{0.5\linewidth}{0.5pt}\end{center}

\subsection{2. Conditional Risk}\label{conditional-risk}

\subsubsection{Definition 2.1: Pointwise Risk}\label{definition-2.1-pointwise-risk}

The pointwise risk of expert \(f_m\) at input \(x\):

\[R_m(x) = \mathbb{E}[\ell(f_m(x), f^*(x)) \mid X = x]\]

where \(\ell: \mathcal{Y} \times \mathcal{Y} \to \mathbb{R}_{\ge 0}\)
is the loss function and \(f^*: \mathcal{X} \to \mathcal{Y}\) is the true labeling function.

\subsubsection{Definition 2.2: State-Conditioned Expert Risk}\label{definition-2.2-state-conditioned-expert-risk}

\[R_m(s) = \mathbb{E}_{x \sim P(\cdot|s)}[\ell(f_m(x), f^*(x))]\]

This is the conditional expectation in the sense of Kolmogorov. Existence is guaranteed by the Radon-Nikodym Theorem:

\[R_m(s) = \frac{1}{P(X \in s)} \int_{X \in s} \ell(f_m(x), f^*(x)) \, dP_X(x)\]

When \(P(X \in s) = 0\), \(R_m(s)\) can be defined as
\(\lim_{\epsilon \to 0^+} R_m(s_\epsilon)\), where \(s_\epsilon\) is the \(\epsilon\)-neighborhood of \(s\).

\subsubsection{Definition 2.3: Conditional Bayes Risk}\label{definition-2.3-conditional-bayes-risk}

\[R^*(s) = \inf_{f \text{ measurable}} \mathbb{E}[\ell(f(X), Y) \mid X \in s]\]

The expert risk \(R_m(s)\) measures expert \(f_m\)'s performance relative to the \textbf{unknown optimal} \(f^*\), rather than relative to \(R^*(s)\).

\subsubsection{Bias-Variance Decomposition}\label{bias-variance-decomposition}

For squared loss \(\ell(\hat{y}, y) = (\hat{y} - y)^2\), let the regression function
\(r(x) = \mathbb{E}[Y \mid X = x]\):

\[R_m(s) = \underbrace{\mathbb{E}[(f_m(X) - r(X))^2 \mid X \in s]}_{\text{approximation error}} + \underbrace{\mathbb{E}[(r(X) - f^*(X))^2 \mid X \in s]}_{\text{target noise}}\]

In the special case \(f^* = r\) we obtain the \textbf{conditional bias-variance decomposition}:

\[R_m(s) = \mathbb{E}[(f_m(X) - r(X))^2 \mid X \in s] + \mathbb{E}[\text{Var}(Y \mid X) \mid X \in s]\]

\begin{center}\rule{0.5\linewidth}{0.5pt}\end{center}

\subsection{3. SCX Reliability}\label{scx-reliability}

\subsubsection{Definition 3.1: SCX Reliability}\label{definition-3.1-scx-reliability}

\[SCX_m(s) = P(\ell(f_m(x), f^*(x)) < \tau \mid x \in s)\]

where \(\tau > 0\) is the reliability threshold. The empirical estimate is:

\[\widehat{SCX}_m(s) = \frac{|\{x_i \in s : \ell(f_m(x_i), y_i) < \tau\}|}{|s|}\]

\subsubsection{Relationship with IRT}\label{relationship-with-irt}

SCX's \(SCX_m(s)\) can be mapped to the Rasch model:

\[\log\left(\frac{SCX_m(s)}{1 - SCX_m(s)}\right) = \theta_s - b_m\]

where \(\theta_s\) is the ``tractability'' parameter of state \(s\) and \(b_m\) is the ability threshold of expert \(m\).

\begin{center}\rule{0.5\linewidth}{0.5pt}\end{center}

\subsection{4. State Data Value}\label{state-data-value}

\subsubsection{Definition 4.1: State Data Value}\label{definition-4.1-state-data-value}

\[V(s) = \bar{r}(s) \cdot \rho(s) \cdot L(s) \cdot [1 - D(s)] \cdot \max_m SCX_m(s)\]

Meaning of each factor:

\begin{longtable}[]{@{}
  >{\raggedright\arraybackslash}p{(\linewidth - 6\tabcolsep) * \real{0.2500}}
  >{\raggedright\arraybackslash}p{(\linewidth - 6\tabcolsep) * \real{0.2500}}
  >{\raggedright\arraybackslash}p{(\linewidth - 6\tabcolsep) * \real{0.2500}}
  >{\raggedright\arraybackslash}p{(\linewidth - 6\tabcolsep) * \real{0.2500}}@{}}
\toprule\noalign{}
\begin{minipage}[b]{\linewidth}\raggedright
Factor
\end{minipage} & \begin{minipage}[b]{\linewidth}\raggedright
Symbol
\end{minipage} & \begin{minipage}[b]{\linewidth}\raggedright
Meaning
\end{minipage} & \begin{minipage}[b]{\linewidth}\raggedright
Range
\end{minipage} \\
\midrule\noalign{}
\endhead
\bottomrule\noalign{}
\endlastfoot
Mean residual & \(\bar{r}(s)\) & Current average prediction error for state \(s\) &
\([0, L_{\max}]\) \\
State probability & \(\rho(s)\) & Occurrence probability \(P(x \in s)\) of state \(s\) &
\([0, 1]\) \\
Learnability & \(L(s)\) & Degree to which performance in state \(s\) can be improved & \([0, 1]\) \\
Diversity penalty & \(1-D(s)\) & Degree of unsaturation of state \(s\) (\(D\) = redundancy) &
\([0, 1]\) \\
Expert coverage & \(\max_m SCX_m(s)\) & Reliability of the best expert on state \(s\) &
\([0, 1]\) \\
\end{longtable}

\subsubsection{Relationship with Shapley Value}\label{relationship-with-shapley-value}

SCX's \(V(s)\) is a \textbf{proxy Shapley value} that decomposes data-point-level Shapley value computation into: first aggregating into states \(s\) (complexity \(O(2^{|S|})\), with \(|S| \ll N\)), then using multiplicative decomposition (product of five interpretable factors).

\begin{center}\rule{0.5\linewidth}{0.5pt}\end{center}

\subsection{5. Noise \& Learnability}\label{noise-learnability}

\subsubsection{Definition 5.1: Noise Score}\label{definition-5.1-noise-score}

\[\text{NoiseScore}(x_i) = r_i \cdot \frac{1}{\rho(s_i) + \varepsilon} \cdot [1 - C(s_i)]\]

where \(C(s)\) is the internal consistency of state \(s\):

\[C(s) = 1 - \frac{1}{|s|} \sum_{x_i \in s} \frac{|\ell(f_m(x_i), y_i) - \bar{\ell}(s)|}{\max \ell - \min \ell}\]

\subsubsection{Definition 5.2: Learnability}\label{definition-5.2-learnability}

\[L(s) = \inf_{f \in \mathcal{F}} \left[ \frac{R(f|s) - R^*(s)}{R^*(s)} \right]^{-1}\]

where
\(R(f|s) = \mathbb{E}[\ell(f(X), Y) | X \in s]\) and \(R^*(s) = \inf_f R(f|s)\)
is the Bayes risk. \(L(s) \to 0\) indicates that the state is unlearnable under the current function class \(\mathcal{F}\).

\subsubsection{Definition 5.3: Redundancy}\label{definition-5.3-redundancy}

State redundancy estimate:

\[D(s) = \rho(s) \cdot \bigl(1 - \bar{r}(s)\bigr) \cdot \text{Sim}(s) \cdot \bigl(1 - \text{Boundary}(s)\bigr)\]

where \(\text{Sim}(s)\) is the intra-state similarity and \(\text{Boundary}(s)\)
is the boundary sample proportion.

\begin{center}\rule{0.5\linewidth}{0.5pt}\end{center}

\subsection{6. Expert Routing}\label{expert-routing}

\subsubsection{Definition 6.1: State-Conditioned Expert Weight}\label{definition-6.1-state-conditioned-expert-weight}

\[w_m(x) \propto \exp\left(-\alpha \cdot \sum_{s \in \mathcal{S}} \gamma_s(x) \cdot \hat{R}_m(s)\right)\]

Normalization ensures \(\sum_{m=1}^M w_m(x) = 1\), with \(\alpha > 0\) being the inverse temperature parameter.

\subsubsection{Definition 6.2: Hard Routing}\label{definition-6.2-hard-routing}

\[m^*(x) = \arg\min_{m \in [M]} \sum_{s \in \mathcal{S}} \gamma_s(x) \cdot R_m(s)\]

\begin{center}\rule{0.5\linewidth}{0.5pt}\end{center}

\subsection{7. Convergence Properties}\label{convergence-properties}

\subsubsection{Theorem 7.1: Concentration of Empirical Risk Estimation}\label{theorem-7.1-concentration-of-empirical-risk-estimation}

Let \(\ell \in [0, L_{\max}]\) and state \(s\) have \(n_s\) samples. For any
\(\delta > 0\), with probability at least \(1 - \delta\):

\[|\hat{R}_m(s) - R_m(s)| \leq L_{\max} \sqrt{\frac{\log(2/\delta)}{2n_s}}\]

\subsubsection{Theorem 7.2: Minimax Lower Bound}\label{theorem-7.2-minimax-lower-bound}

There exists a distribution such that any estimator \(\hat{R}_m(s)\) satisfies:

\[\mathbb{E}[|\hat{R}_m(s) - R_m(s)|] \geq \frac{\sigma(s)}{2\sqrt{n_s}}\]

where \(\sigma^2(s) = \text{Var}(\ell(f_m(X), f^*(X)) \mid X \in s)\).

\begin{center}\rule{0.5\linewidth}{0.5pt}\end{center}

\subsection{Notation Table}\label{notation-table}

\begin{longtable}[]{@{}lll@{}}
\toprule\noalign{}
Symbol & Meaning & Type \\
\midrule\noalign{}
\endhead
\bottomrule\noalign{}
\endlastfoot
\(\mathcal{X}\) & Input space & Metric space \\
\(\mathcal{Y}\) & Label space & \(\mathbb{R}^d\) \\
\(f^*: \mathcal{X} \to \mathcal{Y}\) & True labeling function & Unknown oracle \\
\(f_m: \mathcal{X} \to \mathcal{Y}\) & The \(m\)-th expert model &
Computable function \\
\(M\) & Total number of experts & \(\mathbb{N}\) \\
\(\mathcal{S}\) & State space & \(|\mathcal{S}| = K\) \\
\(s\) & Single state & Measurable subset of \(\mathcal{X}\) \\
\(\gamma_s(x)\) & Soft state assignment & \([0,1]\) \\
\(\ell\) & Loss function & \(\mathbb{R}_{\ge 0}\) \\
\(\tau\) & Reliability threshold & \(\mathbb{R}_{>0}\) \\
\(R_m(s)\) & State-conditioned expert risk & \(\mathbb{R}_{\ge 0}\) \\
\(SCX_m(s)\) & SCX reliability metric & \([0, 1]\) \\
\(V(s)\) & State data value & \(\mathbb{R}_{\ge 0}\) \\
\(\bar{r}(s)\) & State mean residual & \([0, L_{\max}]\) \\
\(\rho(s)\) & State probability & \([0, 1]\) \\
\(L(s)\) & Learnability & \([0, 1]\) \\
\(C(s)\) & Internal consistency & \([0, 1]\) \\
\(N(s)\) & Noise proportion & \([0, 1]\) \\
\(D(s)\) & Redundancy & \([0, 1]\) \\
\(\phi(x)\) & Embedding mapping & \(\mathbb{R}^d\) \\
\(\alpha\) & Inverse temperature parameter & \(\mathbb{R}_{>0}\) \\
\(\eta(x)\) & Regression function \(\mathbb{E}[Y \mid X=x]\) & \(\mathbb{R}\) \\
\(R^*(s)\) & Conditional Bayes risk & \(\mathbb{R}_{\ge 0}\) \\
\end{longtable}

\begin{center}\rule{0.5\linewidth}{0.5pt}\end{center}

\subsection{References}\label{references}

\begin{enumerate}
\def\labelenumi{\arabic{enumi}.}
\tightlist
\item
  Kolmogorov, A. N. (1933). \emph{Grundbegriffe der Wahrscheinlichkeitsrechnung}.
\item
  Wald, A. (1950). \emph{Statistical Decision Functions}.
\item
  Shapley, L. S. (1953). A value for n-person games. \emph{Contributions to the Theory of Games}.
\item
  Tishby, N., Pereira, F. C., \& Bialek, W. (1999). The information bottleneck method. \emph{Allerton Conference}.
\item
  SCX Core Framework Mathematical Analysis. \texttt{01\_SCX\_Core\_Framework\_Math\_Analysis.md}
\item
  SCX Mathematical Foundations and Proof Construction. \texttt{05\_Math\_Roots\_and\_Proofs.md}
\end{enumerate}

\end{document}
